% !TEX program = xelatex
\documentclass[12pt,a4paper,openany]{book}
\usepackage[margin=2.5cm,headheight=15pt]{geometry}
\usepackage{fontspec}
\usepackage{xeCJK}
\setmainfont{Times New Roman}
\setCJKmainfont{Noto Sans TC}
\setCJKsansfont{Noto Sans TC}
\setCJKmonofont{Noto Sans TC}
\usepackage{graphicx}
\usepackage{booktabs,tabularx,array}
\usepackage{xcolor}
\usepackage{listings}
\usepackage{tikz}
\usetikzlibrary{arrows.meta,positioning,shapes.geometric}
\usepackage{hyperref}
\hypersetup{colorlinks=true,linkcolor=teal!60!black,urlcolor=blue!60!black,citecolor=teal!60!black,
  pdftitle={LINE Messaging API 註冊與應用教學},pdfauthor={黃笙誠}}
\usepackage{fancyhdr}
\pagestyle{fancy}
\fancyhf{}
\fancyhead[LE,RO]{LINE Messaging API 教學}
\fancyhead[RE,LO]{黃笙誠｜413630764}
\fancyfoot[C]{\thepage}
\setlength{\parindent}{2em}
\setlength{\parskip}{0.35em}
\renewcommand{\contentsname}{目錄}
\renewcommand{\listfigurename}{圖目錄}
\renewcommand{\listtablename}{表目錄}
\renewcommand{\figurename}{圖}
\renewcommand{\tablename}{表}
\renewcommand{\chaptername}{第}
\renewcommand{\lstlistingname}{程式碼}
\renewcommand{\lstlistlistingname}{程式碼目錄}
\definecolor{codebg}{RGB}{247,249,252}
\definecolor{codeblue}{RGB}{30,75,130}
\lstdefinestyle{pythonstyle}{
  language=Python,basicstyle=\ttfamily\small,keywordstyle=\color{codeblue}\bfseries,
  commentstyle=\color{green!40!black},stringstyle=\color{red!55!black},
  numbers=left,numberstyle=\tiny\color{gray},numbersep=9pt,
  frame=single,rulecolor=\color{black!18},backgroundcolor=\color{codebg},
  breaklines=true,showstringspaces=false,tabsize=4,
  captionpos=b,columns=fullflexible
}
\lstset{style=pythonstyle}

\begin{document}
\frontmatter
\begin{titlepage}
  \centering
  \vspace*{2cm}
  {\large 專業排版與設計｜115 學年度期中作品\par}
  \vspace{2.2cm}
  {\Huge\bfseries LINE Messaging API\par}
  \vspace{0.4cm}
  {\LARGE 註冊申請、Webhook 設定與機器人實作\par}
  \vspace{1.8cm}
  \begin{tikzpicture}[node distance=1.1cm, every node/.style={draw=teal!60!black,rounded corners,fill=teal!5,minimum height=0.85cm,minimum width=2.7cm,font=\small},>=Latex]
    \node (u) {LINE 使用者};
    \node[right=of u] (l) {LINE 平台};
    \node[right=of l] (b) {Flask Bot};
    \draw[->,thick] (u) -- node[above,draw=none,fill=white,minimum width=0cm] {訊息} (l);
    \draw[->,thick] (l) -- node[above,draw=none,fill=white,minimum width=0cm] {Webhook} (b);
    \draw[<-,thick] (l.south) -- ++(0,-0.75) -| node[pos=0.25,below,draw=none,fill=white,minimum width=0cm] {Reply API} (b.south);
  \end{tikzpicture}
  \vfill
  \begin{tabular}{rl}
    班級：& 資管 3A\\
    座號：& 39\\
    學號：& 413630764\\
    姓名：& 黃笙誠\\
    出席次數：& 5\\
    日期：& 2026 年 11 月 13 日
  \end{tabular}
  \vspace{1cm}
\end{titlepage}

\chapter*{摘要}
\addcontentsline{toc}{chapter}{摘要}
本文件以繁體中文說明 LINE Messaging API 的概念、LINE Developers 帳號與 Provider／Channel 建立流程、憑證管理、Webhook 公開測試，以及 Python Flask 回聲機器人的實作。讀者可依章節完成從建立官方帳號到驗證訊息往返的基本流程，並了解部署時的安全與除錯重點。範例採用 LINE Bot SDK v3；實際管理介面與 API 欄位可能隨 LINE 官方更新，操作時應以官方文件為準。本文不包含任何真實憑證。

\tableofcontents
\listoffigures
\listoftables
\lstlistoflistings

\mainmatter
\chapter{Messaging API 與整體流程}
\section{Messaging API 是什麼}
LINE Messaging API 是讓開發者透過 LINE 官方帳號與使用者進行程式化訊息互動的服務。使用者在 LINE 傳送訊息後，LINE 平台會將事件送至開發者設定的 Webhook；後端驗證事件並呼叫 Messaging API，便可回覆訊息或執行其他應用邏輯。它不是使用者帳號的自動化工具，而是以官方帳號、Channel 權限及平台規範為基礎的服務。

常見用途包括商店營業資訊與訂單查詢、課程通知、客服分流、預約提醒、活動報名及內部工作流程通知。簡單的關鍵字回覆適合入門；需要個人化或跨系統資料時，則可在 Webhook 後端加入資料庫、身分驗證和商業邏輯。

\section{角色與訊息往返}
Provider 是 LINE Developers Console 中管理服務的組織或專案容器；Messaging API Channel 則代表一個 bot 與其設定、權限和憑證。官方帳號提供使用者可加入的入口。Webhook 是由 LINE 發出的 HTTPS POST 事件通知；Reply API 則由後端主動向 LINE 平台提出回覆要求。這些角色分離可讓憑證與功能設定更容易管理。

\begin{figure}[htbp]
\centering
\begin{tikzpicture}[node distance=1.2cm,box/.style={draw=blue!55!black,rounded corners,fill=blue!4,align=center,minimum width=3.0cm,minimum height=0.9cm},>=Latex]
\node[box] (user) {使用者\\傳送文字};
\node[box,right=of user] (line) {LINE 平台\\建立事件並簽章};
\node[box,right=of line] (server) {Flask 伺服器\\驗證及處理};
\draw[->,thick] (user)--node[above,draw=none,fill=white]{訊息}(line);
\draw[->,thick] (line)--node[above,draw=none,fill=white]{HTTPS POST}(server);
\draw[->,thick] (server.south) -- ++(0,-0.65) -| node[pos=0.25,below,draw=none,fill=white]{replyToken + 回覆}(line.south);
\draw[->,thick] (line.north) -- ++(0,0.65) -| node[pos=0.75,above,draw=none,fill=white]{顯示回覆}(user.north);
\end{tikzpicture}
\caption{LINE 訊息事件與回覆的基本資料流}
\label{fig:flow}
\end{figure}

\section{開發前準備}
開始前需有可登入 LINE Developers Console 的 LINE 帳號、可管理的 LINE 官方帳號，以及可執行 Python 的電腦。正式 Webhook 需能由 LINE 平台透過 HTTPS 存取；本機開發可用 ngrok 等 HTTPS 通道暫時轉送。測試帳號應先加入官方帳號好友，且需確認官方帳號的回應設定不會與程式回覆互相干擾。

\chapter{註冊申請與 Channel 建立}
\section{登入與建立 Provider}
前往 LINE Developers Console，使用 LINE 帳號登入並完成必要的服務條款與開發者資料設定。第一次使用時，建立 Provider 並以容易辨識的專案或單位名稱命名。Provider 用於集中管理相關 Channels；命名應避免放入密碼、token 或個人敏感資料。若為團隊專案，應由組織管理者建立並以正式成員權限管理，而非共用個人帳密。

\section{建立 Messaging API Channel}
在 Provider 中新增 Channel，選取 Messaging API，依表單填入應用名稱、簡介、類別、聯絡信箱及必要資訊，閱讀並接受當時介面所示條款。建立後進入 Channel 的 Messaging API 設定頁，確認其關聯的官方帳號。LINE Console 介面名稱可能調整，應確認所選類型確實為 Messaging API，而不是 LINE Login 等其他 Channel 類型。

\begin{table}[htbp]
\centering
\caption{註冊階段與檢查重點}
\label{tab:setup}
\begin{tabularx}{\textwidth}{>{\raggedright\arraybackslash}p{2.7cm}XX}
\toprule
階段 & 操作 & 檢查重點 \\
\midrule
登入 & 登入 LINE Developers Console & 使用可持續管理服務的帳號 \\
Provider & 新增 Provider & 名稱可辨識，不含機密資訊 \\
Channel & 建立 Messaging API Channel & 確認 Channel 類型與基本資料 \\
官方帳號 & 依 Console 指示建立或連結 & 確認擁有管理權與測試入口 \\
Messaging API & 開啟設定頁檢查憑證及 Webhook & 憑證只存後端，不放在前端 \\
\bottomrule
\end{tabularx}
\end{table}

\section{官方帳號與管理權限}
Messaging API Channel 需與官方帳號配合，使用者才能找到並加入 bot。依 Console 流程完成官方帳號建立或連結，再到官方帳號管理介面確認個人檔案、好友加入方式及基本回應設定。團隊協作時，透過官方角色或管理權限邀請成員，遵循最小權限原則；離職或專案結束時應移除不再需要的權限。

\chapter{憑證、環境與 Webhook}
\section{Channel Secret 與 Access Token}
Channel Secret 用於驗證 Webhook 簽章，確認請求由持有該 Channel Secret 的 LINE 平台所簽署；Channel Access Token 則用於後端呼叫 Messaging API。兩者用途不同，不可互換。依官方管理介面取得 Secret，並透過 Issue 或 Reissue 產生 Access Token（介面名稱依當時版本為準）。若 token 遺失或疑似外洩，應立即重新發行並更新部署環境。

憑證必須放在伺服器環境變數或受控祕密管理服務，不可寫進程式碼、前端、公開儲存庫或截圖。此專案提供 `.env.example` 作為變數名稱範本；實際 `.env` 已列入 `.gitignore`。分享錯誤訊息前也要確認沒有印出授權標頭、完整請求內容或憑證。

\section{Webhook URL 與簽章驗證}
Messaging API 設定頁填入公開的 HTTPS Webhook URL，例如 `https://你的網域/callback`，啟用 Use webhook，再使用 Verify 測試連線。若服務在本機的 8000 port，可用 ngrok 將 HTTPS 網址轉送至本機；ngrok 顯示的暫時網址改變時，也必須更新 Console 的 Webhook URL。正式部署應使用穩定網域、TLS 憑證及可用的服務監控。

LINE 送來的 `X-Line-Signature` 是以 Channel Secret 對原始請求本文進行簽章後的值。應直接讓官方 SDK 驗證原始本文與簽章，不要先解析、重排 JSON 再驗證。簽章無效時回傳錯誤並停止處理，不能相信未驗證的事件內容。Webhook 端點也應快速回應，耗時工作可移交佇列處理，避免平台重送或逾時。

\section{測試檢查清單}
\begin{enumerate}
  \item 安裝套件，將 `.env.example` 複製成 `.env`，填入測試 Channel 憑證。
  \item 啟動 Flask，先以瀏覽器檢查 `/health` 是否回傳狀態正常。
  \item 建立 HTTPS 通道，將公開網址的 `/callback` 填入 Webhook URL。
  \item 在 Console 啟用 Use webhook 並執行 Verify；確認伺服器收到 POST。
  \item 由測試 LINE 帳號傳送文字，檢查事件、簽章驗證與回覆結果。
  \item 測試無效簽章應被拒絕；檢查日誌不可包含 token 或 secret。
\end{enumerate}

\chapter{Python 回聲機器人實作}
\section{環境安裝與設定}
範例使用 Python、Flask、python-dotenv 與 LINE Bot SDK v3。建議在虛擬環境安裝 `requirements.txt` 中的套件。以 `.env.example` 複製出本機 `.env`，填入兩個 Channel 憑證。程式啟動時會檢查必要變數；若未設定會明確停止，避免服務表面上啟動卻無法處理事件。請勿把 `.env` 加入 Git。

\section{程式流程}
\texttt{/callback} 接收 LINE 的 POST，讀取 \texttt{X-Line-Signature} 和原始本文，再交給 \texttt{WebhookHandler} 驗證。驗證通過後，SDK 依事件類型呼叫 \texttt{handle\_text\_message}；處理器以事件中的 \texttt{reply\_token} 呼叫 Reply API，將使用者文字原樣回傳。非文字事件不會匹配此處理器。回覆 token 有有效期限及使用次數限制；若要稍後或主動通知使用者，需依 LINE 官方規則選擇其他訊息 API 與合適的 token。

\section{回聲機器人原始碼}
下列程式以行號呈現完整範例。行號由 `listings` 套件自動產生，實際執行版本與專案中的 `app.py` 一致。
\lstinputlisting[caption={Python Flask LINE 回聲機器人},label={lst:app},language=Python]{app.py}

\section{常見問題與除錯}
若 Verify 失敗，先確認伺服器可從公網存取、HTTPS 憑證有效、路徑正確，並檢查 Flask 路由是否接受 POST。若收到 400，檢查 Channel Secret 是否屬於同一個 Channel，且代理伺服器沒有改寫請求本文。若 API 回覆失敗，確認 Access Token 未過期、官方帳號與 Channel 關聯正確，並確認 replyToken 尚未使用。不要為了除錯停用簽章驗證；可記錄事件類型、HTTP 狀態和去識別化追蹤資訊。

\begin{table}[htbp]
\centering
\caption{常見故障及初步排查方式}
\label{tab:troubleshoot}
\begin{tabularx}{\textwidth}{>{\raggedright\arraybackslash}p{3.1cm}XX}
\toprule
現象 & 可能原因 & 建議檢查 \\
\midrule
Webhook Verify 失敗 & 網址不可達、路徑或 TLS 錯誤 & HTTPS、轉送服務、路由與伺服器日誌 \\
簽章驗證失敗 & Secret 不符或本文被改寫 & Channel 對應、原始本文與代理設定 \\
回覆 API 錯誤 & Token 無效或 replyToken 失效 & Token 狀態、事件時序及 API 回應碼 \\
本機啟動即停止 & 環境變數缺漏 & `.env` 檔名、欄位與啟動工作目錄 \\
\bottomrule
\end{tabularx}
\end{table}

\chapter{延伸應用、安全與部署}
\section{可延伸的互動功能}
完成文字回聲後，可依服務需求加入關鍵字選單、Quick Reply、Flex Message、Rich Menu、圖片或位置事件。Push Message 適用於符合平台條件的主動通知，但須遵循訊息費用、使用者同意及官方政策。對接訂單或學生資料時，應先設計使用者身分識別、資料最小化與授權邏輯，不應只憑使用者傳入的文字操作敏感資料。

\section{安全與可靠性}
正式部署不得使用 Flask 開發伺服器；應採用正式 WSGI 伺服器、TLS 與受控祕密管理。限制管理介面權限，定期輪替憑證，依需求設定請求大小限制、速率限制及錯誤監控。記錄資訊以除錯所需為限，避免保存不必要的訊息內容、個資與簽章。Webhook 處理應具備冪等性設計，以應對重送事件；重要事件可記錄事件識別資訊並避免重複執行交易。

\section{版本控管與交付}
以 Git 追蹤原始碼、文件和設定範本，排除 `.env`、虛擬環境與快取。每次提交聚焦一項可理解的變更，並在提交前檢查差異是否含憑證。若要備份到 GitHub，先在使用者已登入且有權限的帳號下建立私人或公開儲存庫，再設定遠端並推送；公開儲存庫尤其不可含任何真實憑證。提交次數本身不代表品質，應確保每次提交具明確內容與說明。

\chapter{結論}
本教學完成 LINE Messaging API 的基本註冊概念、Provider 與 Channel 設定、憑證管理、Webhook 公開測試及可執行的 Flask 回聲機器人。核心流程是：使用者事件經 HTTPS Webhook 送至後端，後端驗證簽章並依事件呼叫合適的 Messaging API。實際上線前仍需依官方最新文件確認介面、權限、訊息限制及政策，並落實憑證保護、TLS、日誌去識別化與錯誤監控。

\backmatter
\chapter{本文件 LaTeX 原始碼}
本章依考題要求，以 \verb|\lstinputlisting| 將本文件完整 LaTeX 原始碼排版並自動列入程式碼目錄。為避免讀者誤將憑證寫入文件，原始碼範例僅包含無祕密的教學內容。
\lstinputlisting[caption={本文件 LaTeX 原始碼：main.tex},label={lst:latex},language={[LaTeX]TeX}]{main.tex}

\end{document}